\chapter{Nyelvi alapok és HTTP-határok}

Az A--Z példában már láttuk a fő építőelemeket működés közben. Ebben a fejezetben ezeket napi webfejlesztési munkafolyamattá rendezzük: a forrásfájlokat explicit modulgráfba szervezzük, a statementeket egyértelműen lezárjuk, typed értékekkel számolunk, a szöveget a határokon normalizáljuk, a domainlogikát handlerekben tartjuk, majd typed route-on és biztonságos renderelésen keresztül lépünk át HTTP/HTML világba.

Hasznos munkaszabály: \emph{a bemenetet a határon parse-old és normalizáld, a domainben typed értékekkel dolgozz, explicit query-ken keresztül perzisztálj, és csak a lehető legkésőbb renderelj}. A fejezet sorrendje ezt követi.

\section{Forrásstruktúra és statement-határok}

\subsection*{Modulok}

A \code{main.rw} a modulgráf belépési pontja:

\begin{lstlisting}[language=RWLang]
mod models;
mod queries;
mod components;
mod pages;
mod actions;
\end{lstlisting}

A V1-ben a \code{mod} application-root relatív modult tölt be, de a deklarációit nem önti a globális névtérbe. Például a \code{models.rw} deklarációi más modulból \code{models::Article} alakban érhetők el. A saját modulon belüli lokális deklaráció rövid névvel is hivatkozható. Nincs általános Rust-szerű \code{pub} export-rendszer.

\subsection*{Modul-névterek és modulok közötti hívások}

A moduldeklaráció egy application-root relatív forrásmodult vesz fel az explicit source graphba; nem másolja annak deklarációit az aktuális scope-ba. A kanonikus leképezés közvetlen:

\begin{lstlisting}
models.rw              -> models
catalog/queries.rw     -> catalog::queries
catalog/pages.rw       -> catalog::pages
\end{lstlisting}

Egy belépési pont például több modult tölthet be:

\begin{lstlisting}[language=RWLang]
mod models;
mod catalog::queries;
mod catalog::pages;

route catalogIndex GET "/catalog"
    => catalog::pages::index;
\end{lstlisting}

Modulhatáron át kvalifikált nevet használj. Például:

\begin{lstlisting}[language=RWLang]
let recent = catalog::queries::recent(db);
// tipus modulhataron at: models::Article
\end{lstlisting}

A \code{catalog/queries.rw} tehát a \code{catalog::queries} névtérhez tartozik. A deklaráció saját modulján belül a rövid lokális név továbbra is használható. Így a helyi kód rövid marad, a forrásmodulok közötti dependency viszont látható.

A modulútvonal mindig az application rootból indul. Nincs \code{mod.rw} fallback, automatikus könyvtár-discovery, \code{../}, \code{./}, \code{self::}, \code{super::}, \code{crate::} vagy wildcard import. Minden forrásmodult explicit deklarálj, még a normál top-level deklarációk előtt.

A nyelvi builtinok -- például \code{trim(...)}, \code{sin(...)} és \code{len(...)} -- nem alkalmazásmodul-deklarációk, ezért továbbra is kvalifikáció nélkül használhatók.

\subsection*{Statement terminátorok}

Az RWLangban az egyszerű statementeket explicit pontosvessző zárja. A sortörés whitespace, nem statement-határ; nincs automatikus semicolon insertion. Ez determinisztikussá teszi a többsoros kifejezéseket és egyértelműbbé a generált, illetve auditált kódot.

\begin{lstlisting}[language=RWLang]
let subtotal = price
    * quantity
    + shipping;
set retries = retries + 1;
authorize article owner authorUsername or role Publisher;
return Ok(json(subtotal));
\end{lstlisting}

A blokkos statementeket és blokkos deklarációkat a kapcsos zárójel zárja, utánuk nincs pontosvessző:

\begin{lstlisting}[language=RWLang]
if ready {
    set state = 1;
}
while state < 3 {
    set state = state + 1;
}
\end{lstlisting}

Ugyanez érvényes tranzakción belül is: az önálló mutating query-hívás és a business-audit statement egyszerű statement, ezért \code{;} zárja. A top-level blokkos deklarációk -- például \code{model}, \code{page fn}, \code{action fn} -- és a control-flow blokkok után nincs \code{;}. A blokk nélküli top-level deklarációk viszont pontosvesszősök: a \code{mod path;} és a \code{route ... => handler;} is explicit \code{;} terminátort kap. Többsoros route esetén is csak ez a záró terminátor fejezi be a deklarációt.

\section{Typed számítás webalkalmazás-kódban}

\subsection*{Kifejezések és aritmetika}

Az RWLang explicit operátorprecedenciát használ. Szorosabbtól lazább felé: unáris \code{!}; szorzás, osztás és maradékképzés; összeadás és kivonás; shiftek; összehasonlítások; bitenkénti AND/XOR/OR; logikai AND; logikai OR. Ha javítja az olvashatóságot, használj zárójelet.

\begin{lstlisting}[language=RWLang]
let total = 12 + 3 * 4;
let remainder = 17 % 5;
let shifted = 3 << 2;
let masked = 12 & 10;
let allowed = ready && authorized;
let fallback = cached || fresh;
let denied = !allowed;
\end{lstlisting}

Az operátorszerződés szigorúan typed:

\begin{center}
\begin{tabularx}{\textwidth}{@{}l l Y@{}}
\toprule
\textbf{Operandusok} & \textbf{Operátorok} & \textbf{Eredmény} \\
\midrule
\code{Int}, azonos típus & \code{+ - * / \%} & checked \code{Int}. \\
\code{F32}, azonos típus & \code{+ - * / \%} & \code{F32}; véges eredmény szükséges. \\
\code{Decimal}, azonos típus & \code{+ - * / \%} & checked \code{Decimal}. \\
\code{Int}, azonos típus & \code{<< >> \& \^{} |} & shift/bitművelet, \code{Int}. \\
\code{Bool}, azonos típus & \code{\&\& ||} & \code{Bool}; short-circuit kiértékelés. \\
\code{Bool} & \code{!} & logikai negáció. \\
\code{String}, azonos típus & \code{+} & budgetelt String-konkatenáció. \\
Numerikus, azonos típus & \code{< <= > >=} & \code{Bool}. \\
Támogatott azonos skalártípus & \code{== !=} & \code{Bool}. \\
\bottomrule
\end{tabularx}
\end{center}

Nincs implicit keverés \code{Int}, \code{F32} és \code{Decimal} között. Lebegőpontos számításhoz az integert explicit \code{toF32(value)} hívással alakítsd át. Overflow, nullával osztás/maradékképzés, hibás shift vagy nem véges \code{F32} eredmény fail-closed hibát ad.

A numerikus builtinok:

\begin{center}
\begin{tabularx}{\textwidth}{@{}l l Y@{}}
\toprule
\textbf{Függvény} & \textbf{Típus} & \textbf{Feladat} \\
\midrule
\code{sin(x)}, \code{cos(x)} & \code{F32 -> F32} & Trigonometrikus függvények. \\
\code{sqrt(x)} & \code{F32 -> F32} & Négyzetgyök; hibás tartomány elutasítva. \\
\code{abs(x)} & \code{Int -> Int} vagy \code{F32 -> F32} & Abszolútérték. \\
\code{ln(x)}, \code{log10(x)} & \code{F32 -> F32} & Természetes/tízes alapú logaritmus. \\
\code{log(x, base)} & \code{F32, F32 -> F32} & Explicit alapú logaritmus. \\
\code{exp(x)} & \code{F32 -> F32} & Exponenciális függvény. \\
\code{pow(x, y)} & \code{F32, F32 -> F32} & Lebegőpontos hatványozás. \\
\code{round/floor/ceil(x)} & \code{F32 -> F32} & Kerekítési műveletek. \\
\code{toF32(x)} & \code{Int -> F32} & Explicit integer--\code{F32} konverzió. \\
\code{monotonicNanos()} & \code{() -> Int} & Process-lokális monoton nanoszekundum-számláló. \\
\bottomrule
\end{tabularx}
\end{center}

Minden \code{F32} eredménynek végesnek kell maradnia. A \code{sqrt(-1.0f32)} és \code{ln(-1.0f32)} jellegű domainhibák NaN/végtelen helyett hibát adnak. A \code{monotonicNanos()} eltelt idő mérésére való, nem falióra-időbélyeg; a tőle függő public-cache kimenetet a compiler elutasítja.

\subsection*{Gyakori webes aritmetikai minták}

Webalkalmazásban ritkán kell mindenhol haladó matematika, néhány numerikus minta viszont folyamatosan visszatér. Ezeket typed értékekkel és látható kóddal kezeld, ne string-konverziókba vagy SQL-részletekbe rejtsd.

Lapozásnál validált integerekből számolj offsetet:

\begin{lstlisting}[language=RWLang]
let pageSize = 25;
let offset = (page - 1) * pageSize;
\end{lstlisting}

Integer flag/mask esetén használhatók a bitműveletek, de üzleti feltételekhez az olvasható boolean logika az alap. A \code{&&} és \code{||} short-circuit módon működik:

\begin{lstlisting}[language=RWLang]
let mayPublish = authenticated
    && ownsArticle
    && !locked;
\end{lstlisting}

Pénzhez \code{Decimal}-t használj. \code{F32}-t akkor válassz, ha maga a feladat lebegőpontos: például geometria, jelfeldolgozás, statisztika vagy explicit matematikai függvény. Pénzösszeget ne alakíts \code{F32}-re csak azért, hogy math builtint lehessen rá alkalmazni.

\subsection*{String- és szövegfüggvények}

A string műveletek typed, Unicode-tudatos builtinok a korlátozott runtime-on belül:

\begin{lstlisting}[language=RWLang]
let cleaned = trim(text);
let left = trimStart(text);
let right = trimEnd(text);
let normalized = lower(cleaned);
let chars = stringLen(cleaned);
let rewritten = replace(normalized, " ", "-");
let part = substring(cleaned, 1, 3);
let first = indexOf(cleaned, "rw");
let last = lastIndexOf(cleaned, "a");
let ch = charAt("RWLang", 1);
let repeated = repeat("rw", 3);
let parts = split(rewritten, "-");
\end{lstlisting}

\begin{center}
\begin{tabularx}{\textwidth}{@{}Y Y@{}}
\toprule
\textbf{Függvény} & \textbf{Szignatúra} \\
\midrule
\code{stringLen(text)} & \code{String -> Int} \\
\code{trim/trimStart/trimEnd(text)} & \code{String -> String} \\
\code{lower/upper(text)} & \code{String -> String} \\
\code{contains(text, needle)} & \code{String, String -> Bool} \\
\code{startsWith/endsWith(text, part)} & \code{String, String -> Bool} \\
\code{replace(text, search, replacement)} & \code{String, String, String -> String} \\
\code{split(text, delimiter)} & \code{String, String -> List<String>} \\
\code{substring(text, start[, length])} & \code{String, Int[, Int] -> String} \\
\code{indexOf/lastIndexOf(text, needle)} & \code{String, String -> Int} \\
\code{charAt(text, index)} & \code{String, Int -> String} \\
\code{repeat(text, count)} & \code{String, Int -> String} \\
\bottomrule
\end{tabularx}
\end{center}

A \code{stringLen}, \code{substring}, \code{indexOf}, \code{lastIndexOf} és \code{charAt} Unicode scalar value pozíciókat használ, nem UTF-8 byte indexet. Sikertelen keresésnél \code{-1} az eredmény; negatív vagy érvénytelen index fail-closed hibát ad. A \code{replace} nem enged üres keresőszöveget, a \code{split} nem enged üres delimitert és legfeljebb 4096 elemet ad, a \code{repeat} pedig az eredmény létrehozása előtt bekerül az allocation budgetbe. A \code{List<String>} typed compute collection marad \code{len(parts)} és ellenőrzött \code{parts[index]} használattal.

Korlátozott reguláris kifejezéses feldolgozáshoz elérhető a \code{regexMatch}, \code{regexReplace} és \code{regexCaptures}; ezek továbbra is a validációs és resource limitek alatt futnak.

\subsection*{Gyakori szövegkezelési minták}

A szövegmanipulációt kezeld boundary-feladatként. Keresési és filter inputot érdemes összehasonlítás előtt normalizálni, de az eredeti értéket tartsd meg, ha a domainben számít a pontos írásmód.

\begin{lstlisting}[language=RWLang]
let queryText = trim(q);
let normalizedQuery = lower(queryText);
let hasRwlang = contains(normalizedQuery, "rwlang");
\end{lstlisting}

A \code{replace} determinisztikus helyettesítésre való, nem általános parser helyett:

\begin{lstlisting}[language=RWLang]
let displayTag = replace(trim(tag), "_", " ");
\end{lstlisting}

Szeletelésnél és pozíciókeresésnél az RWLang Unicode scalar-value pozíciókat használ UTF-8 byte offset helyett. Így a \code{substring}, \code{charAt} és indexfüggvények alkalmazásszövegnél nem teszik ki a fejlesztőt encoding-részleteknek. Strukturált azonosítóknál -- URL, slug, email, dátum, UUID -- viszont a dedikált RWLang típusokat és validátorokat használd, ne string függvényekből építsd újra a szabályokat.

\subsection*{Melyik eszközt válaszd alkalmazáskódban?}

\begin{center}
\begin{tabularx}{\textwidth}{@{}l Y@{}}
\toprule
\textbf{Webfejlesztési feladat} & \textbf{Preferált RWLang megoldás} \\
\midrule
Keresési/filter szöveg normalizálása & \code{trim}, majd \code{lower}/\code{upper} csak akkor, ha az összehasonlítás case-insensitive. \\
Ismert token helyettesítése & \code{replace}; regexet csak valódi mintakereséshez használj. \\
Email, URL, slug, dátum, UUID & Dedikált typed érték és validáció, ne ad-hoc string parsing. \\
Pénz/összegzés & \code{Decimal}. \\
Lapozás/számlálók & \code{Int} checked aritmetikával. \\
Üzleti feltételek & \code{Bool} short-circuit \code{&&}, \code{||}, \code{!} operátorokkal. \\
Tudományos/statisztikai matematika & \code{F32} és explicit math builtin; előbb döntsd el, hogy a lebegőpont valóban illik-e a domainhez. \\
Eltelt idő mérése & \code{monotonicNanos()}, soha ne falióra-időbélyegként. \\
\bottomrule
\end{tabularx}
\end{center}

\section{Domainértékek és handlerkód}

\subsection*{Modellek és üzleti típusok}

\begin{lstlisting}[language=RWLang]
model Article {
    id: Int
    slug: Slug
    title: String
    body: String
}
\end{lstlisting}

A gyakori típusok: \code{String}, \code{Int}, \code{Bool}, továbbá \code{Date}, \code{DateTime}, \code{Uuid}, \code{Decimal} és \code{Slug}. Pénzhez \code{Decimal}-t használj, ne lebegőpontos számot.

\subsection*{Page és action}

A page GET-szerű megjelenítési logika:

\begin{lstlisting}[language=RWLang]
page fn show(ctx: PageContext, db: Db, id: Int)
    -> Result<Html, PageError> {
    let article = articleById(db, id)?;
    return Ok(html { <h1>{{ article.title }}</h1> });
}
\end{lstlisting}

Az action állapotváltoztató kérés kezelője:

\begin{lstlisting}[language=RWLang]
action fn create(ctx: ActionContext, db: Db, title: String)
    -> Result<Redirect, PageError> {
    transaction db {
        insertArticle(tx, title)?;
    }
    return Ok(redirect("/admin/articles"));
}
\end{lstlisting}

\subsection*{Változók és hibaterjedés}

A \code{let} typed lokális értéket hoz létre. Általános \code{mut} nincs a V1-ben; compute kódban explicit \code{set name = expr} újraértékadás használható az eredeti statikus típus megtartásával és instruction budget alatt. A \code{?} a typed Result hibáját továbbítja.

\subsection*{Enum}

Zárt választási halmazhoz ne szabad szöveget használj, hanem enumot. Ez form-, route-, JSON- és domain-contractként is erősebb.

\subsection*{Object}

Összetartozó model/query/page műveletek domain object alatt csoportosíthatók. Ez különösen nagyobb alkalmazásban segít a domain API körülhatárolásában.


\section{HTTP- és prezentációs határ}

\subsection*{Route}

\begin{lstlisting}[language=RWLang]
route article GET "/cikk/:slug<Slug>" => article;
\end{lstlisting}

A path paraméter típusa a route-ban deklarált. Hibás input nem jut el a handlerig typed értékként.

\subsection*{HTML escape}

\begin{lstlisting}[language=RWLang]
return Ok(html {
    <h1>{{ article.title }}</h1>
});
\end{lstlisting}

A \code{{ ... }} HTML-escaped. Adatbázisból érkező string sem kivétel. Raw HTML API nincs.

\subsection*{Lista és optional érték}

A listát iterálni kell, az optional modelt pedig explicit jelenlétvizsgálat után lehet mezőzni. A compiler nem engedi, hogy \code{List<Model>} vagy ellenőrizetlen \code{Model?} úgy viselkedjen, mintha egyetlen rekord lenne.

\begin{lstlisting}[language=RWLang]
@for article in articles {
    <li>{{ article.title }}</li>
}

@if article {
    <h1>{{ article.title }}</h1>
}
\end{lstlisting}

A query-cardinality pontos futásidejű szerződését -- például hogy a \code{Model?} több sor esetén hibát jelent -- az adatbázis-fejezet tárgyalja.

\subsection*{Typed URL}

\begin{lstlisting}[language=RWLang]
<a @href(article, article.slug)>Megnyitas</a>
<form method="post" @action(deleteArticle, article.id)>
\end{lstlisting}

\subsection*{Preferáld a route helpereket}

Saját alkalmazás-route-ok URL-jét ne másold kézzel több helyre. HTML-ben POST formhoz az \code{@action(...)} és linkhez az \code{@href(...)} a preferált minta, mert a route neve és typed paraméterei maradnak a kapcsolat forrásai. A beépített runtime endpointok -- például a \code{/__rw/auth/logout} -- természetesen nem alkalmazás-route helperek.


Ne építs dinamikus \code{href="{{ value }}"} linket. A route helper ismeri a célt és a paramétertípust.

\subsection*{Component és layout}

\begin{lstlisting}[language=RWLang]
layout fn Main(title: String) -> Html {
    html {
        <html>
            <head><title>{{ title }}</title></head>
            <body><main>@content</main></body>
        </html>
    }
}
\end{lstlisting}

Page-ben:

\begin{lstlisting}[language=RWLang]
return Ok(html {
    @layout(Main, article.title) {
        <article><h1>{{ article.title }}</h1></article>
    }
});
\end{lstlisting}

A layout pontosan egy \code{@content} slotot tartalmazhat; component/layout csak explicit typed paramétereket lát.
